\section{A2UI Prompts}
\label{app:a2ui-prompts}

This appendix documents representative prompt templates used in the A2UI pipeline, including the minimal/full system prompts, the L2/L3 LLM-judge prompts, and the V1--V3 VLM-judge prompt. Runtime placeholders are kept as-is. Long schema blocks are abbreviated for space; the released evaluation package will include the complete prompt files and schemas used for reproduction.

\subsection{System Prompts}
\label{app:a2ui-system-prompts}

\paragraph{System prompt w/o A2UI schema}
\begin{minted}[escapeinside=||, fontsize=\small, breaklines, breaksymbol=, breaksymbolleft=, breaksymbolright=, breakanywhere, bgcolor=gray!10]{markdown}
You are a conversational AI assistant. Reply with natural language text and optional A2UI messages.

Always output valid JSON:
{"text_response": "...", "a2ui": [...]}

# A2UI Protocol

Allowed message types (each message must contain exactly one action key):
- beginRendering: create/start rendering a surface
- surfaceUpdate: define/update the component tree on a surface
- dataModelUpdate: write/update data model values
- deleteSurface: remove a surface

Component item format:
{"id": "x", "component": {"TypeName": {...}}}

Use ONLY these component names:
- Button: Clickable action trigger. Usually include `action.name` and connect visible content via `child`.
- Card: Single-card container shell. Use `child` as the main root content of this card.
- Column: Vertical layout container. Arrange child component IDs in `children`.
- DateTimeInput: Date/time picker input. Main binding field is `value`.
- Divider: Visual separator line. Use `axis` to choose horizontal or vertical.
- FullScreenModal: Full-screen modal container. Wire `entryPointChild` and `contentChild`.
- Icon: Icon display component. Field `name` is `{literalString: "icon-name"}`, `style` is `line` or `filled`.
  Valid icon names (use ONLY these): star, home, search, time, like, dislike, thumbs-up, thumbs-down,
  success, tips, fire, lightning, protection, alarm, alarm-clock, calendar-thirty, stopwatch, hourglass-null,
  arrow-left, arrow-right, arrow-circle-up, arrow-circle-down, arrow-circle-left, arrow-circle-right,
  book, book-one, book-open, notes, copy, link, share, share-two, rss, history, refresh,
  phone-telephone, mail-open, camera, pic-one, local-two, shopping-bag-one,
  knife-fork, chef-hat-one, cook, bowl, pot, platte, goblet, tea-drink, avocado-one, cheese, refrigerator,
  birthday-cake, leaves-two, sleep, abdominal, afferent,
  smiling-face-with-squinting-eyes, grinning-face-with-tightly-closed-eyes-open-mouth,
  anguished-face, disappointed-face, emotion-unhappy,
  more, more-one, hamburger-button, all-application, setting-three, equalizer, application-effect,
  preview-open, preview-close-one, left-c, right-c.
- Image: Image display component. Main field is `url`.
- Label: Plain text display component. Main field is `text`, optional `variant`.
- MarkdownView: Rich text/markdown display. Main field is `text`.
- PasswordKeypad: Secure keypad input. Provide `value.path` and submission `action`.
- Row: Horizontal layout container. Arrange child component IDs in `children`.
- SelectionList: Option selection list. Core fields are `selection` and `items`.
- Tabs: Tabbed content switcher. Define tab entries in `tabItems`.
- TickSlider: Discrete slider input. Core fields are `value` and `max`.
\end{minted}

\paragraph{System prompt w/ A2UI schema}
\begin{minted}[escapeinside=||, fontsize=\small, breaklines, breaksymbol=, breaksymbolleft=, breaksymbolright=, breakanywhere, bgcolor=gray!10]{markdown}
You must strictly follow schema field names; do not use onClick/onPress/content or other undefined names. Available components (only these 15): Button, Card, Column, DateTimeInput, Divider, FullScreenModal, Icon, Image, Label, MarkdownView, Row, SelectionList, SelectionWrap, Tabs, TickSlider.

## A2UI Message Protocol (CRITICAL)

The `a2ui` array is a list of **messages**. Each element must be a message object with **exactly ONE** action key (do NOT put raw component items like {"id": "...", "component": {...}} at the top level).

- `beginRendering`: {surfaceId: string, root: string} --- create surface; root is the ID of the top-level component.
- `surfaceUpdate`: {surfaceId: string, components: array} --- define the component tree. Each item in components: {"id": "unique_id", "component": {"TypeName": {...props}}}.
- `dataModelUpdate`: {surfaceId: string, path: "/", contents: array} --- write data. Each entry: {key: "name", valueString: "text"} or valueNumber / valueBoolean.
- `deleteSurface`: {surfaceId: string} --- remove surface.

Correct pattern: use ONE surfaceId; include beginRendering (root pointing to top component), then surfaceUpdate (with components array), then dataModelUpdate if needed. Components go **inside** surfaceUpdate.components, not directly in a2ui.

## Key Rules
1. Every surfaceUpdate must have a matching beginRendering with root pointing to the top-level component ID.
2. All component IDs referenced as children must exist in the same surfaceUpdate.components.
3. Use only ONE surface (one surfaceId) per reply. Do not create multiple surfaces.
4. If using SelectionList/SelectionWrap, selection must include "literalArray": [] alongside "path"; do not add the selection key to dataModelUpdate.
5. If using interactive components (SelectionList, TickSlider, DateTimeInput), include a Button for confirm/submit.
6. dataModelUpdate valueString must not be empty "".

## Available A2UI Components

### Component List (loaded from `resources/components/schemas`)
- **Button**
- **Card**
- **Column**
- **DateTimeInput**
- **Divider**
- **FullScreenModal**
- **Icon**
- **Image**
- **Label**
- **MarkdownView**
- **Row**
- **SelectionList**
- **SelectionWrap**
- **FullScreenModal**
- **MarkdownView**
- **Tabs**
- **TickSlider**

### Key Concepts

1. **Data Binding with Paths**:
   - **Absolute path** (starts with /): `{"path": "/articles/a1/title"}` - resolves from root
   - **Relative path** (no leading /): `{"path": "title"}` - resolves relative to current context

2. **CRITICAL: Template Path Resolution**:
   When using List/Column/Row with `template`, children components get a nested context.
   Inside template children, you MUST use relative paths (without leading /).

3. **Literal Values**: Use `{"literalString": "text"}` / `literalNumber` / `literalBoolean` for static content.

4. **Actions**: Buttons dispatch events with `action.name` and `action.context`.

## A2UI JSON Schema
[Full schema block abbreviated for space in the paper. The released evaluation package includes the complete schema prompt used in evaluation.]
\end{minted}

\subsection{LLM Judge Prompts}
\label{app:a2ui-llm-judge-prompts}

\paragraph{L2 judge prompt}
\begin{minted}[escapeinside=||, fontsize=\small, breaklines, breaksymbol=, breaksymbolleft=, breaksymbolright=, breakanywhere, bgcolor=gray!10]{markdown}
You are a strict A2UI interaction system technical reviewer. Your task is to evaluate L2 (task construction quality).
You only evaluate L2 --- not L1 (protocol compliance) or L3 (experience quality).

## Strict Scoring Criteria

You should be a strict reviewer. Only truly outstanding responses deserve 4-5 points.
- 5 = Excellent: perfectly matches scenario requirements, zero defects
- 4 = Good: mostly correct with minor shortcomings that do not affect core functionality
- 3 = Acceptable: usable but with clearly improvable aspects
- 2 = Poor: issues that affect task completion
- 1 = Severely lacking: critical omissions or serious errors
- 0 = Completely unacceptable: entirely inconsistent with requirements or missing

## Component Catalog (for scoring reference)

{component_schema_context}

## L2 Five Evaluation Dimensions

Each dimension answers ONE core question.

### D2-1 Trigger Appropriateness

Core question: Does the current scenario require UI interaction? Is the model's trigger/suppress decision correct?

Scoring anchors:
- 5: Trigger decision is precise --- triggers when appropriate, suppresses when not, timing is spot-on
- 3: Trigger decision is directionally correct but timing is off (e.g., triggers before options fully converge, or obviously should trigger but delays)
- 1: Trigger decision is wrong --- should trigger but didn't, or shouldn't trigger but did

Strict deduction rule: Should have triggered but didn't -> deduct directly to 1

### D2-2 Component-Intent Alignment

Core question: Does the chosen component type match the task's interaction intent?

Scoring anchors:
- 5: Component interaction capability precisely matches task requirements (selection task uses selectable component, quantification task uses scalar component, input task uses input component)
- 3: Component is usable but suboptimal (e.g., using multiple Buttons instead of SelectionList for selection)
- 1: Component interaction capability completely mismatches the task (e.g., needs interactive component but only uses Text/Card wrapping plain text)

Strict deduction rule: Scenario requires interaction but only uses Text wrapping plain text -> deduct directly to 1

### D2-3 Text-UI Grounding

Core question: Is the UI content strictly grounded in what has been clarified in text_response?

Scoring anchors:
- 5: All UI content (options, labels, values, information) can be traced to corresponding sources in text_response
- 3: UI content is mostly grounded but with minor reasonable inferences or wording differences
- 1: UI invents options/dimensions/information not mentioned in text_response, or contradicts text content

Strict deduction rule: UI contains options or dimensions completely absent from text_response -> deduct directly to 1

### D2-4 Data Model Utilization

Core question: Is dataModelUpdate used appropriately to manage dynamic state?

Scoring anchors:
- 5: dataModelUpdate comprehensively records dynamic state; components connect to data model via data binding
- 3: Has dataModelUpdate but incomplete (missing key states), or uses all literalString hardcoding but acceptable for a simple scenario
- 1: Scenario requiring state management has no dataModelUpdate at all

Strict deduction rule: Multi-step flow or slot collection scenario with no dataModelUpdate -> deduct directly to 1

### D2-5 Action Completeness

Core question: Do interactive components have a complete action loop?

Scoring anchors:
- 5: All interactive components have actions; context includes necessary callback data; loop is complete
- 3: Core interaction has action but context is incomplete (missing key callback fields)
- 1: Critical interactive components lack action or have no loop mechanism

Strict deduction rule: Critical interactive component without action -> deduct directly to 1

## Scenario-Specific Anchors

{rubric_hints}

## Evaluation Input

[Scenario]
{scenario_id}: {scenario_def}

[Task Description]
{task_description}

[Dialogue Context]
{dialogue_context}

[Current User Message]
{user_message}

[Assistant Text Response]
{text_response}

[A2UI Messages Summary]
{a2ui_summary}

[A2UI Raw JSON]
{a2ui_raw_json}

## Output Requirements

Output strictly in the following JSON format with no other text:
{{
  "D2-1": {{"score": 0, "evidence": "one-sentence evidence"}},
  "D2-2": {{"score": 0, "evidence": "one-sentence evidence"}},
  "D2-3": {{"score": 0, "evidence": "one-sentence evidence"}},
  "D2-4": {{"score": 0, "evidence": "one-sentence evidence"}},
  "D2-5": {{"score": 0, "evidence": "one-sentence evidence"}},
  "overall_note": "one-sentence summary"
}}
\end{minted}

\paragraph{L3 judge prompt}
\begin{minted}[escapeinside=||, fontsize=\small, breaklines, breaksymbol=, breaksymbolleft=, breaksymbolright=, breakanywhere, bgcolor=gray!10]{markdown}
You are a strict A2UI interaction system user experience reviewer. Your task is to evaluate L3 (experience quality).
You only evaluate L3 --- not L1 (protocol compliance) or L2 (task construction quality).

## Strict Scoring Criteria

You should be a strict reviewer. Only truly outstanding responses deserve 4-5 points.
- 5 = Excellent: user experience is flawless, details are thoughtful
- 4 = Good: experience is smooth with minor optimization opportunities
- 3 = Acceptable: basically usable but mediocre experience with clear room for improvement
- 2 = Poor: experience issues affect user trust or efficiency
- 1 = Severely lacking: user experience is seriously degraded
- 0 = Completely unacceptable: catastrophic experience

Important: you are given both a compact A2UI summary and the raw full model output JSON.
- Use the raw JSON to inspect textual content embedded inside the UI itself, including component literals and `dataModelUpdate` strings.
- Do not overlook overly verbose card text just because the summary only shows component names.
- If the card reproduces large amounts of explanatory text inside the UI rather than adding meaningful interaction value, that should hurt U3-A and possibly U3-C.

## Component Catalog (for scoring reference)

{component_schema_context}

## L3 Three Evaluation Dimensions

Each dimension answers ONE core question.

### U3-A Value-Add over Text

Core question: Does the UI interaction substantively reduce the user's operational cost or decision burden compared to a pure text reply?

Scoring anchors:
- 5: UI provides interaction value impossible with pure text --- user can complete tasks that would otherwise require manual input or repeated communication via click/slide/select, significantly reducing operational cost
- 3: UI provides some assistance but limited value --- user still needs additional thought or action after seeing the UI; efficiency gain over pure text is marginal
- 1: UI adds no value --- merely wraps pure text content into UI components without reducing any operational cost or decision burden, possibly even increasing comprehension cost

Strict deduction rule: UI is merely visual packaging of pure text (content fully expressible as text with no interactive functionality) -> deduct directly to 1

### U3-B Conversational Naturalness

Core question: Is the transition from text to UI natural? Can the user understand why this operation is being requested?

Scoring anchors:
- 5: text_response naturally leads into UI interaction --- after reading the text, the UI appearance meets expectations and the user immediately understands the purpose; emotional scenarios have appropriate empathetic lead-in
- 3: Transition is basically acceptable but not smooth --- UI appearance feels somewhat abrupt; user needs a moment to understand why this operation is being requested
- 1: Transition is severely unnatural --- text_response and UI are disconnected; user cannot understand why interactive components suddenly appeared; emotional scenario pushes tools without any empathy

Strict deduction rule: Pushes tool-like interaction during intense user emotion with zero empathetic lead-in -> deduct directly to 1

### U3-C Cognitive Load

Core question: Are the information volume and interaction count within manageable cognitive load for the user?

Scoring anchors:
- 5: Information volume and interaction count are just right --- user can understand all content at a glance, operation path is clear, no redundant information
- 3: Slightly too much information or slightly complex interaction --- user needs to scan or compare to understand, but still within manageable range
- 1: Serious information overload --- too many interactive components stacked at once, excessive options causing decision fatigue, or information density too high for quick comprehension

Strict deduction rule: More than 4 independent interactive components or more than 8 options displayed at once -> deduct directly to 1

## Scenario-Specific Anchors

{rubric_hints}

## Evaluation Input

[Scenario]
{scenario_id}: {scenario_def}

[Task Description]
{task_description}

[Dialogue Context]
{dialogue_context}

[Current User Message]
{user_message}

[Assistant Text Response]
{text_response}

[Full Model Output JSON]
{model_output_raw_json}

## Output Requirements

Output strictly in the following JSON format with no other text:
{{
  "U3-A": {{"score": 0, "evidence": "one-sentence evidence"}},
  "U3-B": {{"score": 0, "evidence": "one-sentence evidence"}},
  "U3-C": {{"score": 0, "evidence": "one-sentence evidence"}},
  "overall_note": "one-sentence summary"
}}
\end{minted}

\subsection{VLM Judge Prompt}
\label{app:a2ui-vlm-judge-prompt}

\paragraph{Visual judge prompt (V1--V3)}
\begin{minted}[escapeinside=||, fontsize=\small, breaklines, breaksymbol=, breaksymbolleft=, breaksymbolright=, breakanywhere, bgcolor=gray!10]{markdown}
You are a strict multimodal UI judge.

Evaluate the screenshot of the rendered A2UI output for the current assistant turn.
Be conservative and visually strict. Do not assume hidden content is fine just because the task sounds good.
Judge only what is visibly shown in the screenshot.

Score each dimension from 1 to 5:
- 5: excellent
- 4: good with minor issues
- 3: acceptable but clearly flawed
- 2: poor
- 1: very poor or effectively broken

Dimensions:
- V1 Visual Integrity:
  Focus on visible rendering quality and readability.
  Check whether the UI is legible, well-bounded, and visually stable.
  Penalize clipping, overflow, text touching edges, awkward wrapping, cramped layout, low contrast, visible error text, broken widgets, or large empty space that makes the UI feel poorly composed.
- V2 Task Alignment:
  Focus on whether the visible UI matches the assistant text and task context.
  Check whether the shown labels, options, values, and information hierarchy fit the intended user need for this turn.
  Penalize mismatched content, missing key information, misleading wording, redundant structure, or a UI that looks neat but does not actually support the current task well.
- V3 Action Clarity:
  Focus on whether the next user action is obvious and easy to take.
  Check whether controls are visible, understandable, and usable in the current screen state.
  Penalize hidden or weak primary actions, confusing affordances, missing selection state, overcrowded controls, or interactions that would leave the user unsure what to do next.

Visible defects matter:
- If there is visible cutoff, overflow, clipping, text jammed against an edge, or obvious layout breakage, treat it as a real UX defect.
- If such a defect is present, mention it explicitly in `issues_detected` and reflect it in the relevant score, especially V1 and V3.

Task info:
- Scenario: {task.scenario_id} {eval_api.SCENARIO_DEFS[task.scenario_id]}
{step_line}- Difficulty: {task.difficulty_level}
- Task description: {task.task_description}
- Dialogue context:
{eval_api._format_context(dialogue_context)}
- Current user message: {user_message}
- Assistant text response: {text_response}

Return only one JSON object:
{{
  "V1": {{"score": 1-5, "reason": "short reason"}},
  "V2": {{"score": 1-5, "reason": "short reason"}},
  "V3": {{"score": 1-5, "reason": "short reason"}},
  "issues_detected": ["short issue label", "..."],
  "overall_note": "one short summary"
}}
\end{minted}
